\section{精确计数的更紧界}\label{app:quantitative}

我们通过收缩布尔张量网络来证明定理~\ref{thm:main}。
这两个运行时间反映了网络中导线的数量。
含有 $m\le Dn$ 个子句的公式给出一个具有 $O_{k,D}(n)$
条导线的网络。对于任意公式，我们使用
第~\ref{sec:hashes} 节的奇偶检验。引理~\ref{lem:total-order}
的次数界使我们能用 $O_k(n\log\log n)$ 条导线来索引其多项式
提升的乘积中保留的一个单项式。我们在证明它们所用的收缩
定理之前给出这两种编码。

\subsection{布尔张量收缩}
\label{subsec:tensor-interface}

布尔张量网络是一个二部多重图 $G$，其中每条边携带一个比特。
每个顶点 $v$ 有一个关联边的有序列表 $\delta(v)$ 和一个张量
$T_v$，后者为这些边的每个布尔赋值指派一个环元素。平行边
具有不同的端口。该网络的\emph{收缩}为
\[
 Z=\sum_{x\in\{0,1\}^{E(G)}}\prod_{v\in V(G)} T_v(x|_{\delta(v)}).
\]
孤立顶点贡献其标量元素。我们用边数来度量网络的
大小。

\begin{theorem}[布尔张量收缩]\label{thm:tensor-contraction}
固定正常数 $C,A,b$。设 $s\ge2$ 为整数，并设
$\mathcal L$ 是一个二部张量网络的图与张量描述的总长度，
该网络至多有 $Cs$ 条边且最大度至多为 $\lceil s^{9/10}\rceil$。
假设顶点数被 $s$ 的某个固定多项式所界，且各元素属于
\[
 \mathcal B_{P_0,\ell}
   =(\mathbb Z/2^{P_0}\mathbb Z)[V]/(V^\ell),
 \qquad 1\le P_0,\ell\le Cs.
\]
假设对每个顶点及其端口的每种布尔取值一致地，其元素可由
该描述在时间 $(s+\mathcal L)^b2^{A s^{9/10}}$ 与空间
$(s+\mathcal L)^b$ 内计算出来。那么该收缩可在
\[
 (s+\mathcal L)^{O(1)}
  2^{O(s/\sqrt{\log s})}
\]
位运算与 $(s+\mathcal L)^{O(1)}$ 比特空间内确定性计算。
常数只依赖于假设中的固定界。特别地，
允许 $\ell=1$。
\end{theorem}

第~\ref{app:tensor} 节中的算法用沿导线的递归取代单次全局
泰勒展开。在每次调用时，它形变导线的等同关系，使得在选定的
插值节点处，只有同一组内的端口发生相互作用。在该节点附近，
组之间的相互作用按称为模式的变量展开。权重为 $j$ 的模式的
系数可被 $2^j$ 整除，因此只有总权重低于工作精度的单项式
才能有贡献。用节点的 2 进赋值渐缩 Hermite 重数，从而
界定该精度。在 $M$ 条导线与 $u$ 个保留的赋值层级下，
模式单项式的数目为 $2^{O(M/\sqrt{u})}$，且每个单项式分解为
至多 $M/8$ 条导线的收缩之积。在递归深度 $d$ 处，
$M=O(s8^{-d})$ 且 $u=\Theta((\log s)/2^d)$。因此这些指数
按几何级数求和为 $O(s/\sqrt{\log s})$。

\subsection{固定子句密度}\label{subsec:sharp-fixed-density}

\begin{proof}[定理~\ref{thm:main} 中固定密度界的证明]
给每个子句一个顶点，该顶点携带其布尔满足表，子句中
每个不同变量对应一个端口。重言式子句可以删除，空子句
使计数为零。对于具有 $d\ge1$ 个端口的变量，要求它们的
比特一致。对 $d=1$ 附加一元张量 $(1,1)$，对 $d=2$ 附加
二元相等张量。对 $d\ge3$，取一棵具有 $d$ 个外部端口、
内部顶点为三价相等张量的树。这样的树有 $d-2$ 个内部
顶点，其端口上的每个公共值唯一地延伸到其所有
内部边。

至多有 $km$ 个子句端口。该图有 $O_k(n+m)$ 条边和顶点，
最大度至多为 $\max\{k,3\}$。用二元相等张量把每条边
细分一次。旧顶点与新顶点构成二部网络的两侧。细分边上
的两个比特必须相等，因此这保持收缩不变。出现变量上的
每个满足赋值贡献一，其他每个赋值贡献零。

取 $s=n+m$，在 $\mathbb Z/2^{n+1}\mathbb Z$ 上应用
定理~\ref{thm:tensor-contraction}，取 $P_0=n+1$ 与 $\ell=1$。
对固定的 $k$，局部表的大小有界。次数与精度的界对
所有充分大的 $s$ 成立，而有界实例可直接枚举。若有 $u$ 个
声明变量未被使用，则把收缩乘以 $2^u$ 并对
$2^{n+1}$ 取模。其最小非负剩余为 $\#F$，因为
$0\le\#F\le2^n$。这包括空公式，其空网络的收缩为一。
时间为 $s^{O_k(1)}2^{O_k(s/\sqrt{\log s})}$，空间为多项式。
当 $m\le Dn$ 时，有 $n\le s\le(1+D)n$，给出指数
$O_{k,D}(n/\sqrt{\log n})$。
\end{proof}

\subsection{任意子句密度}\label{subsec:sharp-dense}

第~\ref{sec:hashes} 节中的稀疏化与贪心行选择把这一情形
归约为计数稀疏多项式方程的公共零点。我们保留
第~\ref{sec:setup} 节的奇偶提升。其次数界使我们能用
$O(n\log\log n)$ 条布尔导线，通过检验—变量关联的
子集来编码一个单项式。

\begin{lemma}[通过收缩计数稀疏多项式系统]
\label{lem:tensor-sparse-query}
固定整数 $k,c\ge1$。设 $P_1,\ldots,P_a$ 是 $\mathbb F_2$ 上
关于 $n\ge20$ 个声明布尔变量的多项式。假设
$a\le cn$，且每个 $P_i$ 被列为至多
$\lceil\log_2 n\rceil^c$ 个次数至多为 $k$ 的单项式。公共零点的
数目可在 $n^{O_{k,c}(1)}2^{O_{k,c}(n\log\log n/\sqrt{\log n})}$
位运算与 $n^{O_{k,c}(1)}$ 比特空间内确定性计算。
\end{lemma}

\begin{proof}
令 $\lambda=\lceil\log_2 n\rceil$ 与
$\nu=\lceil\log_2\lambda\rceil$。
在每个单项式内，把变量的正幂次替换为其一次幂，并按模二
合并重复的单项式。检验 $i$ 的作用域由其剩余单项式中的
变量组成。其大小至多为 $k\lambda^c$。
设 $N_0$ 为各作用域大小之和，其中若变量 $v$ 属于检验 $i$，
则对 $(i,v)$ 只计一次。于是 $N_0\le ak\lambda^c$。
以所有可能的方式固定属于多于
\[
 D_*=\max\{1,\lceil N_0/n\rceil\}\lceil\sqrt{\lambda}\rceil
\]
个作用域的变量。这样的变量至多有 $n/\sqrt{\lambda}$ 个。它们的
赋值把原计数划分为至多 $2^{n/\sqrt{\lambda}}$ 个分支计数。
在每个分支上代入固定比特并再次合并重复单项式。检验 $i$ 的
自由作用域 $I_i$ 的大小仍至多为 $k\lambda^c$，且每个自由
变量出现在至多 $D_*=\lambda^{O_{k,c}(1)}$ 个作用域中。不出现
在任何作用域中的自由变量仍保持声明状态。常值一方程使
分支计数为零。我们可以丢弃常值零方程；若没有剩余方程，
则该分支的答案为二的自由于变量数目的幂。下文 $a$ 表示
剩余方程的数目，$n$ 表示原始变量数。

\paragraph{关联上的多项式。}
给每个满足 $v\in I_i$ 的对 $(i,v)$ 一个自己的未定元
$z_{iv}$。记 $z_i=(z_{iv})_{v\in I_i}$。用整数系数求规范化
分支多项式剩余单项式之和，以 $z_{iv}$ 代替 $x_v$，构成
$\widetilde L_i(z_i)$。我们使用环
\[
 \mathcal B=(\mathbb Z/2^p\mathbb Z)[V]/(V^\ell),
 \qquad p=\ell=n+1 .
\]
取 $h(Y)=Y(Y-1)$ 以及引理~\ref{lem:rational-idempotent} 中的
整系数多项式 $e_j$。将其系数模 $2^p$ 约化，并构成因子
\[
 B_i(z_i,V)=\sum_{j=0}^{\ell-1}
   e_j(\widetilde L_i(z_i))\,h(\widetilde L_i(z_i))^j V^j .
\]
设 $\overline B_i$ 为其在 $z_i$ 中的多线性代表，它通过在 $V$ 的
每个系数中把每个关联未定元的正幂次替换为一次幂而得到。它在
$z_i$ 的每个布尔赋值上与 $B_i$ 一致。在 $\mathcal B[z]$ 中定义
$F(z,V)=\prod_i\overline B_i(z_i,V)$。不同检验的关联未定元
互不相交，因此即使在把同一个布尔 $x_v$ 代入变量 $v$ 的所有
副本之前，$F$ 也是多线性的。

$\overline B_i$ 中 $V^j$ 的系数的关联次数至多为 $k(1+2j)$，
因为 $\deg e_j\le1$、$\deg h=2$，且多线性化只会降低单项式的
次数。$\mathcal B$ 中的乘法只保留满足 $\sum_i j_i\le\ell-1$
的项。从而 $F$ 中每个关联单项式的次数至多为
\begin{equation}\label{eq:sharp-incidence-degree}
 K=k\bigl(a+2(\ell-1)\bigr)=O_{k,c}(n).
\end{equation}
设 $\mathcal S_\ell$ 在 $\mathcal B$ 中求 $1,V,\ldots,V^{\ell-1}$
的系数之和。在 $z_{iv}=x_v$ 处，多线性化因子
$\overline B_i$ 与规范化分支多项式的 $B_i$ 一致。
整数 $\widetilde L_i((x_v)_{v\in I_i})$ 与原 $P_i$ 在 $x$ 与
固定比特处取值具有相同的奇偶性。对这些整数应用
引理~\ref{lem:total-order} 表明，分支计数模 $2^p$ 同余于
\begin{equation}\label{eq:sharp-branch-residue}
 \mathcal S_\ell\left(\sum_x F\bigl((z_{iv}=x_v)_{iv},V\bigr)\right),
\end{equation}
其中 $x$ 遍历所有自由布尔赋值。我们在截断环中乘起各检验
因子之后应用系数求和映射。

\paragraph{索引单项式。}
分别把检验与自由变量划分为 $J=\lceil n^{1/4}\rceil$ 组，
每侧各组的大小至多相差一。用 $\mathcal T_g$ 表示检验组，
$\mathcal X_h$ 表示变量组。一组包含 $O_c(n^{3/4})$ 个对象，
因此其关联数在任一侧至多为
$n^{3/4}\lambda^{O_{k,c}(1)}$。设 $N_{gh}$ 为从 $\mathcal T_g$ 到
$\mathcal X_h$ 的关联数，并设 $N=\sum_{g,h}N_{gh}\le N_0$。

对关联单项式，设 $d_{gh}$ 为其连接组 $g,h$ 的入选关联数。
枚举满足 $0\le d_{gh}\le N_{gh}$ 与 $\sum_{g,h}d_{gh}\le K$ 的
整数数组 $(d_{gh})$。至多有
\[
 (K+1)^{J^2}=2^{O_{k,c}(\sqrt n\log n)}
\]
个数组，可用 $J^2$ 个计数器枚举。对固定数组，这一对存在
\[
 C_{gh}=\binom{N_{gh}}{d_{gh}}
\]
个指定大小的子集。在两个组顶点之间引
$\lceil\log_2 C_{gh}\rceil$ 条布尔导线，以字典序索引它们。
当 $C_{gh}=1$ 时，两个端点使用唯一的子集（可能是整个关联集）
且不使用导线。任何未使用的比特串给出张量元素零。为解码
一个有效索引，按顺序取各关联，并用二项式系数计数包含或
不包含每个关联的可能补全。这些具有多项式比特长度，因此
解码使用多项式时间与空间。

所有对上的子集之并指定了 $N$ 个未定元中一个次数至多 $K$ 的
无平方单项式。计入该次数界的所有单项式，由隔板法给出上界
$\binom{N+K}{K}$。因此该数组对应网络中的导线数 $M$ 满足
\begin{align}
 M &\le J^2+\log_2\prod_{g,h} C_{gh}
      \le J^2+\log_2\binom{N+K}{K} \notag\\
   &\le J^2+K\log_2\!\left(\frac{\mathrm e(N+K)}{K}\right)
      = O_{k,c}(n\nu).
 \label{eq:sharp-wire-budget}
\end{align}
这里 $K=\Theta_{k,c}(n)$ 且 $N\le n\lambda^{O_{k,c}(1)}$。
又 $\lceil\log_2 C_{gh}\rceil\le N_{gh}$，故每个顶点的度至多为
其关联数 $n^{3/4}\lambda^{O_{k,c}(1)}$。

\paragraph{张量元素与收缩。}
导线比特的一个有效赋值在每个检验处确定一个子集
$S_i\subseteq I_i$。记 $z_{i,S_i}=\prod_{v\in S_i}z_{iv}$。检验组
$g$ 处的元素是其所指单项式在
$\prod_{i\in\mathcal T_g}\overline B_i(z_i,V)$ 中的系数。由于各
$z_i$ 互不相交，这等于
\[
 \prod_{i\in\mathcal T_g}[z_{i,S_i}]\overline B_i(z_i,V)
 \qquad\text{在 }\mathcal B\text{ 中}.
\]
对变量组 $h$，对 $v\in\mathcal X_h$ 令 $b_v=|\{i:v\in S_i\}|$。
其元素为
\[
 \sum_{x\in\{0,1\}^{\mathcal X_h}}
        \prod_{v\in\mathcal X_h}x_v^{b_v}
   =2^{|\{v\in\mathcal X_h:b_v=0\}|}\pmod{2^p}.
\]
确实，正指数要求 $x_v=1$，而零指数使 $x_v$ 自由。这包括
所有检验作用域之外的变量。

$F$ 的每个单项式恰好出现在一个数组中，并在每个组对处恰有
一个有效索引。它在 $F$ 中的系数是各检验组系数之积。在把具有
相同 $v$ 的关联设为相等之后，其在布尔赋值上的和是各变量组
元素之积。由 \eqref{eq:sharp-incidence-degree}，对这些数组
求和穷尽了所有单项式。从而 \eqref{eq:sharp-branch-residue}
中 $\mathcal S_\ell$ 内 $\mathcal B$ 的元素是这些二部网络的
收缩之和。

为计算一个检验侧系数，设 $\mathbf1_T$ 为 $I_i$ 上支撑集为
$T\subseteq S_i$ 的布尔向量。在 $\mathbf1_T$ 处，该多线性
多项式的值等于其以下标为 $T$ 的子集所索引的系数之和，且
在该处与 $B_i$ 一致。反解此子集和给出
\[
 [z_{i,S_i}]\overline B_i(z_i,V)=
 \sum_{T\subseteq S_i}(-1)^{|S_i|-|T|}B_i(\mathbf1_T,V)
 \qquad\text{在 }\mathcal B\text{ 中}.
\]
每一项花费多项式多个位运算。由所列单项式计算
$\widetilde L_i(\mathbf1_T)$，并使用引理~\ref{lem:rational-idempotent}
中的多项式 $e_j$，它们可在多项式时间内构造。存储
$B_i(\mathbf1_T,V)$ 的 $\ell$ 个系数模 $2^p$。至多有
$2^{|S_i|}\le 2^{k\lambda^c}$ 个子集需要处理。每次累加一个
检验，并在组内于 $\mathcal B$ 中乘起所得系数。这以
$2^{\lambda^{O_{k,c}(1)}}$ 位运算与多项式空间计算一个检验组
元素。变量侧元素与无效导线索引的检查花费多项式时间。
图、次数数组与元素计算过程的描述具有多项式大小。

以规模参数 $s=n\nu$、根精度 $p=n+1$ 与截断阶 $\ell=n+1$
应用定理~\ref{thm:tensor-contraction}。由
\eqref{eq:sharp-wire-budget}，该图有 $2J$ 个顶点与
$O_{k,c}(s)$ 条导线。对所有充分大的 $n$，其最大度至多为
$\lceil s^{9/10}\rceil$，其精度与截断阶为 $O(s)$，其元素计算
时间至多为 $2^{O_{k,c}(s^{9/10})}$。每次收缩随后花费
$n^{O_{k,c}(1)}2^{O_{k,c}(s/\sqrt{\log s})}$ 位运算与多项式
空间。这里 $s/\sqrt{\log s}=O(n\nu/\sqrt\lambda)$。
高关联分支与次数数组至多把运行时间乘以
\[
 2^{\,n/\sqrt{\lambda}+O_{k,c}(\sqrt n\log n)}.
\]
该指数中的两项均为 $O_{k,c}(n\nu/\sqrt{\lambda})$。
顺序处理各分支与数组。在每个分支上，于 $\mathcal B$ 中
求和各收缩并应用 $\mathcal S_\ell$。模 $2^{n+1}$ 相加各分支
剩余即得原计数，因为它介于 $0$ 与 $2^n$ 之间。这证明了
所声称的界，其中有界例外规模采用直接枚举。
\end{proof}

\begin{proof}[定理~\ref{thm:main} 中任意密度界的证明]
如第~\ref{subsec:general-reduction} 节那样，在其原始 $n$ 个
声明变量上规范化该公式。空子句与不含子句的公式有直接
答案。对 $n\ge20$，取
$\sigma=1/\lceil\sqrt{\lambda}\rceil$ 应用引理~\ref{lem:sparsify}，
其中 $\lambda=\lceil\log_2 n\rceil$。该树有
$n^{O_k(1)}2^{O(n/\sqrt{\lambda})}$ 个节点与多项式深度。
每个叶子的出现次数界为
$B=O_k((\sigma^{-1}\log(2/\sigma))^{k-1})=\lambda^{O_k(1)}$，
且至多有 $Bn$ 个非空子句。

对于具有 $m_G$ 个子句的非平凡叶子，使用
引理~\ref{lem:sparse-checks} 的初始奇偶检验与
引理~\ref{lem:greedy-hashes} 的贪心过程。每个初始行的
支撑大小至多为 $D_0$，每个附加行的列表长度为 $P$。
两个界均为 $\lambda^{O_k(1)}$。由 \eqref{eq:hash-query}，每行
给出一个次数至多 $k$、至多含 $2^k\max\{P,D_0\}$ 个单项式的
方程，且每个被查询的系统至多有 $6n$ 个方程。对充分大的
$n$，单个常数 $c=c(k)$ 同时界定方程与变量之比以及单项式
界中 $\lambda$ 的指数。在每个贪心步骤，使用
引理~\ref{lem:tensor-sparse-query} 对每个候选行 $h_\gamma$ 计算
$N(A;h_\gamma)$，并追加取得最小值的一行。该求值器也适用于
被丢弃的候选，其系统仍接受非满足赋值。最后一次调用给出
所选矩阵的 $N(A)$，由引理~\ref{lem:greedy-hashes} 它等于
叶子计数。

构造初始检验与枚举候选（除求值之外）花费
$2^{\lambda^{O_k(1)}}$ 位运算。每个叶子至多有
$3n(2m_G)^P+1=2^{\lambda^{O_k(1)}}$ 次求值。由于 $\lambda$ 的
每个固定幂次都是 $o(n\log\log n/\sqrt{\log n})$，这些开销
与叶子数目被所声称的时间界
\[
 (n+m)^{O_k(1)}2^{O_k(n\log\log n/\sqrt{\log n})}
\]
吸收。这里 $m$ 计输入子句数。深度优先遍历该树，每次处理
一个查询。各叶子的解集互不相交，因此其精确计数用一个
$(n+1)$ 比特的累加器相加。结合每个查询与每条树路径的
多项式空间，这给出所断言的多项式空间界。
\end{proof}

单个算法同时达到这两个界。取
$\lambda=\lceil\log_2 n\rceil$ 与 $\nu=\lceil\log_2\lambda\rceil$，
当 $m\le n\nu$ 时使用子句网络，否则使用奇偶查询。在第一个
分支中 $n\le s=n+m=O(n\log\log n)$，因此即使在增长的密度下
其指数也是 $O_k(n\log\log n/\sqrt{\log n})$。对每个固定 $D$ 与
所有充分大的 $n$，每个满足 $m\le Dn$ 的输入都走该分支且
$s=O_D(n)$。
